% =========================================================
%  Section : La forme quadratique E_8
% =========================================================

\section{La forme quadratique \texorpdfstring{$E_8$}{E8} : un exemple fondamental}

\subsection{Définition et premières propriétés}

Parmi les formes bilinéaires symétriques entières unimodulaires, la forme $E_8$ occupe une
place centrale. Elle est définie sur $\mathbb{Z}^8$ par la matrice
\[
  E_8 =
  \begin{pmatrix}
     2 & -1 &  0 &  0 &  0 &  0 &  0 &  0 \\
    -1 &  2 & -1 &  0 &  0 &  0 &  0 &  0 \\
     0 & -1 &  2 & -1 &  0 &  0 &  0 & -1 \\
     0 &  0 & -1 &  2 & -1 &  0 &  0 &  0 \\
     0 &  0 &  0 & -1 &  2 & -1 &  0 &  0 \\
     0 &  0 &  0 &  0 & -1 &  2 & -1 &  0 \\
     0 &  0 &  0 &  0 &  0 & -1 &  2 &  0 \\
     0 &  0 & -1 &  0 &  0 &  0 &  0 &  2 \\
  \end{pmatrix}.
\]

\begin{proposition}\label{prop:e8_proprietes}
  La forme $Q_{E_8}$ satisfait les propriétés suivantes~:
  \begin{enumerate}
    \item \textbf{Définie positive.} Pour tout $v \in \mathbb{Z}^8 \setminus \{0\}$,
          $Q_{E_8}(v,v) > 0$. La signature est $(8,0)$.
    \item \textbf{Unimodulaire.} $\det E_8 = 1$.
    \item \textbf{Paire.} Pour tout $v \in \mathbb{Z}^8$, $Q_{E_8}(v,v) \equiv 0 \pmod{2}$.
    \item \textbf{Rang minimal.} $E_8$ est l'unique forme définie positive, paire et
          unimodulaire de rang minimal~; ce rang vaut $8$.
  \end{enumerate}
\end{proposition}

\begin{proof}
  \textit{(i)} Le critère de Sylvester s'applique~: les mineurs principaux successifs de
  $E_8$ valent $2, 3, 4, 5, 6, 4, 2, 3$, tous strictement positifs.

  \textit{(ii)} Un développement de Laplace exploitant la structure tridiagonale par blocs
  donne $\det A_{E_8} = 1$.

  \textit{(iii)} Pour tout $v = (v_1, \ldots, v_8) \in \mathbb{Z}^8$,
  \[
    Q_{E_8}(v,v) = 2\sum_{i=1}^8 v_i^2
    - 2\bigl(v_1 v_2 + v_2 v_3 + v_3 v_4 + v_4 v_5 + v_5 v_6 + v_6 v_7 + v_3 v_8\bigr),
  \]
  qui est bien pair pour tout $v \in \mathbb{Z}^8$.

  \textit{(iv)} \textit{L'unicité nécessite un arsenal plus développé.}
\end{proof}

% ---------------------------------------------------------
\subsection{Non-diagonalisabilité sur \texorpdfstring{$\mathbb{Z}$}{Z}}

\begin{proposition}\label{prop:e8_non_diag}
  La forme $Q_{E_8}$ n'est pas isomorphe, sur $\mathbb{Z}$, à une somme directe de formes
  de rang $1$.
\end{proposition}

\begin{proof}
  Comme $Q_{E_8}$ est unimodulaire et définie positive, une diagonalisation éventuelle serait
  nécessairement de la forme $\langle 1 \rangle^{\oplus 8}$. Or $\langle 1 \rangle^{\oplus 8}$
  est impaire — chaque vecteur de base $e_i$ vérifie $Q(e_i, e_i) = 1$ qui est impair —
  tandis que $Q_{E_8}$ est paire. Deux formes de parités différentes ne peuvent être
  isomorphes.
\end{proof}

\begin{remark}
  Sur $\mathbb{Q}$, en revanche, toute forme définie positive est diagonalisable, et l'on a
  $Q_{E_8} \otimes \mathbb{Q} \cong \langle 1 \rangle^{\oplus 8}$. La non-diagonalisabilité
  est donc un phénomène proprement arithmétique, qui ne se voit pas au niveau rationnel.
\end{remark}

% ---------------------------------------------------------
\subsection{Réalisation comme forme d'intersection}

Rappelons que la forme d'intersection d'une $4$-variété compacte orientée fermée simplement
connexe est unimodulaire, et paire si et seulement si la variété est spin. La forme $E_8$
vérifiant ces deux conditions, on peut se demander si elle est réalisable géométriquement.

\medskip

La réponse est positive dans la catégorie topologique~: par le théorème de Freedman, il existe
une $4$-variété topologique compacte fermée orientée simplement connexe $M_{E_8}$, unique à
homéomorphisme près, dont la forme d'intersection est $Q_{E_8}$.

\medskip

En revanche, la forme $E_8$ ne peut pas être réalisée par une variété lisse. C'est une
conséquence directe du théorème de Donaldson~: toute $4$-variété lisse compacte fermée
simplement connexe à forme d'intersection définie positive a une forme diagonalisable sur
$\mathbb{Z}$. Or $Q_{E_8}$ est définie positive et non diagonalisable sur $\mathbb{Z}$
(Proposition~\ref{prop:e8_non_diag}), d'où~:

\begin{corollary}\label{cor:e8_non_lissable}
  La variété $M_{E_8}$ n'admet aucune structure différentielle.
\end{corollary}
